\documentclass[10pt,a4paper]{amsart}
\usepackage{amsmath,amssymb,amsthm,graphicx}
\usepackage{hyperref}
\usepackage{tikz}
%\usepackage{xcolor}
\usepackage{pict2e}
\let\arc\relax
%\usepackage{harpoon}
\usepackage{curves}
%\usepackage{chngcntr}
%\usepackage{apptools}
\usepackage{enumerate}

%\usepackage{verbatim}
%\usepackage[nospace,noadjust]{cite}
\usepackage{accents}
\usetikzlibrary{calc,decorations.markings}
\unitlength1cm

%%%%%%%%%%%%%%FONTS%%%%%%%%%%%%%%%%%%%%%%%%%%%%%%
\newcommand{\R}{\mathbb{R}}
\newcommand{\Z}{\mathbb{Z}}
\newcommand{\N}{\mathbb{N}}
\newcommand{\C}{\mathbb{C}}
\newcommand{\Id}{I}

%%%%%%%%%%%%%%MATHCAL%%%%%%%%%%%%%%%%%%%%%%%
\newcommand{\cK}{\mathcal{K}}
\newcommand{\cE}{\mathcal{E}}
\newcommand{\cR}{\mathcal{R}}
\newcommand{\cW}{\mathcal{W}}

%%%%%%%%%%%%%%%%%%ABBRS%%%%%%%%%%%%%%%%%%%%%%%%%%%%%

%%%%%%% OBJECTS SPECIFIC TO THIS PAPER%%%%%%%%%

\newcommand{\hw}{\widehat w}
\newcommand{\tw}{\widetilde w}

\newcommand{\bQ}{ Q}
\newcommand{\hQ}{{\widehat Q}}
\newcommand{\tQ}{{\widetilde Q}}

\newcommand{\sbQ}{Q^\star}
\newcommand{\shQ}{{\widehat Q^\star}}
\newcommand{\stQ}{{\widetilde Q^\star}}

\newcommand{\smu}{\mu^\star}
\newcommand{\shmu}{\widehat \mu^\star}
\newcommand{\stmu}{\widetilde \mu^\star}

\newcommand{\bmu}{{\mu}}
\newcommand{\hmu}{{\widehat \mu}}
\newcommand{\tmu}{{\widetilde \mu}}


\newcommand{\bP}{P}
\newcommand{\hP}{{\widehat P}}
\newcommand{\tP}{{\widetilde P}}

\newcommand{\sbP}{P^\star}
\newcommand{\shP}{{\widehat P^\star}}
\newcommand{\stP}{{\widetilde P^\star}}

\newcommand{\hR}{{\widehat R}}

\newcommand{\bv}{v}
\newcommand{\hv}{{\widehat v}}
\newcommand{\tv}{{\widetilde v}}

\newcommand{\sbv}{v^\star}
\newcommand{\shv}{{\widehat v^\star}}
\newcommand{\stv}{{\widetilde v^\star}}

\newcommand{\stu}{{\widetilde u^\star}}

\newcommand{\hF}{\widehat F}
\newcommand{\hW}{\widehat W}

\newcommand{\Fw}{\frac{F^2}{w}}
\newcommand{\hFw}{\frac{\hF^2}{\hw}}

\newcommand{\Cau}{\mathcal C_{\Sigma}}
\newcommand{\Cauv}{\mathcal C_{v}}
\newcommand{\Cauhv}{\mathcal C_{\hv}}
\newcommand{\Cautv}{\mathcal C_{\tv}}

\newcommand{\sCau}{\mathcal C_{\Sigma^\star}}
\newcommand{\sCauv}{\mathcal C_{\bv^\star}}
\newcommand{\sCauhv}{\mathcal C_{\hv^\star}}
\newcommand{\sCautv}{\mathcal C_{\stv}}

\newcommand{\Ltwo}{L^2(\Sigma)}
\newcommand{\sLtwo}{L^2(\Sigma^\star)}

\newcommand{\Sigl}{\Sigma^\ell}
\newcommand{\Sigr}{\Sigma^r}

%%%%%%%%%%%%%%%%OTHER SHORTCUTS%%%%%%%%%%%
\newcommand{\nn}{\nonumber}
\newcommand{\ul}{\underline}
\newcommand{\ol}{\overline}
\newcommand{\pa}{\partial}
\newcommand{\ti}{\tilde}
\newcommand{\wti}{\widetilde}
\newcommand{\wha}{\widehat}
\newcommand{\norm}[1]{\lVert#1 \rVert}
\newcommand{\abs}[1]{\lvert#1 \rvert}
%\newcommand{\Id}{\mathbb{I}}
\newcommand{\I}{\mathrm{i}}
\newcommand{\E}{\mathrm{e}}
\newcommand{\re}{\mathop{\mathrm{Re}}}
\newcommand{\im}{\mathop{\mathrm{Im}}}
\newcommand{\eps}{\varepsilon}
\newcommand{\dist}{\mathop{\mathrm{dist}}}

\newtheorem{Theorem}{Theorem}[section]
\newtheorem{Corollary}[Theorem]{Corollary}
\newtheorem{Lemma}[Theorem]{Lemma}
\newtheorem{Proposition}[Theorem]{Proposition}
 { \theoremstyle{definition}
\newtheorem{Definition}[Theorem]{Definition}
\newtheorem{Note}[Theorem]{Note}
\newtheorem{Example}[Theorem]{Example}
\newtheorem{Remark}[Theorem]{Remark} }

\numberwithin{equation}{section}
\usepackage{cleveref}
\crefname{originalequation}{original equation}{original equations}
\Crefname{originalequation}{Original equation}{Original equations}
\makeatletter
\newlabel{ParaConv}{{internal}{source}{}{sourcepointer-1}{}}
\newlabel{ParaConv@cref}{{[originalequation][0][]internal}{[1][0][]source}{}{}{}}
\makeatother
\crefname{originallemma}{original lemma}{original lemmas}
\Crefname{originallemma}{Original lemma}{Original lemmas}
\makeatletter
\newlabel{PropP}{{internal}{source}{}{sourcepointer-2}{}}
\newlabel{PropP@cref}{{[originallemma][0][]internal}{[1][0][]source}{}{}{}}
\makeatother
\title{Four estimates for the norms of the jump matrix differences}
\author{Deift, Percy, Piorkowski, Mateusz}
\date{doi:10.3842/SIGMA.2024.004; CC BY 4.0}
\begin{document}
\maketitle
\setcounter{section}{7}
\setcounter{Theorem}{0}
\noindent\textit{Editorial note.} All delivered mathematical text is the author's own, copied from the pinned SIGMA article source; no statement or proof step was rewritten, reordered or removed. The counted statement is the labeled proposition environment at source lines 1790-1798 and the counted proof is the proof environment at source lines 1799-1898.\par\medskip
\noindent\textit{Editorial note.} The proof states that it follows for the most part a published proposition of the cited literature, the difference being the contribution of the boundary of the disc; that reference is declared as a standard prerequisite of this card and the adaptation carried out in the delivered text is preserved verbatim.\par\medskip
\noindent\textit{Editorial note.} The proof cites published references of the author's bibliography by key; the cited entries are copied line for line from the author's own bibliography into a retained bibliography fragment, and hyperref is added because the publisher class, not the author preamble, supplied the link command.\par\medskip
\noindent\textit{Editorial note.} The retained preamble of this unit was generated mechanically for this delivery: it copies the author's own package, numbering, macro and theorem declaration block from the pinned source (source lines 9-133), with the publisher class and page setup lines removed.\par\medskip
\noindent\textit{Allowed standard prerequisites (editorial).}\begin{itemize}
\item The corresponding estimates for the problem without the boundary contribution. published statement cited by the author at proof line 1800. Supplies the argument for the bulk part of the estimates, which the author states is taken over with the only difference being the boundary contribution.
\item The matching relation between the parametrices and the outer parametrix. same article, the display at source line 1377. Controls the difference between the two parametrices on the two parts of the contour, which is the quantity whose norm is estimated.
\item The properties of the local parametrices near the endpoint. same article, the lemma at source lines 1363-1390. Provides the growth bounds that are used when the estimates are split into the two parts of the contour.
\end{itemize}\medskip
\ifdefined\hypertarget\hypertarget{sourcepointer-1}{}\fi
\noindent\textit{Source reference.} Original equation internal: sources/sigma-2024-004-recurrence-coefficients-for-orthogonal-polyn/pinned-original-source.tex, source lines 1377--1379. The counted statement or proof refers to that the author's displayed matching relation between the parametrices and the outer parametrix; it is the author's own earlier display, definition or statement in the same article and the whole argument of the counted span is retained verbatim.\par\medskip
\ifdefined\hypertarget\hypertarget{sourcepointer-2}{}\fi
\noindent\textit{Source reference.} Original lemma internal: sources/sigma-2024-004-recurrence-coefficients-for-orthogonal-polyn/pinned-original-source.tex, source lines 1363--1390. The counted statement or proof refers to that the author's lemma collecting the properties of the three local parametrices; it is the author's own earlier display, definition or statement in the same article and the whole argument of the counted span is retained verbatim.\par\medskip
\begin{Proposition}\label{PropFourInt}
The following estimates hold for $n \to \infty$:
\begin{enumerate}[$(i)$]\itemsep=0pt
 \item $\Vert \shmu_n (\sbv_n - \shv_n) \Vert_{L^2(\Sigma^\star)} = O\big(\frac{1}{n^{1/2}\log^2 n}\big)$,
 \item $\Vert \smu_n - \shmu_n \Vert_{L^2(\Sigma^\star)} = O\big(\frac{1}{n^{1/2}\log^2 n}\big)$,
 \item $\Vert \shmu_n (\sbv_n - \shv_n) - \shmu_{n+1} (\sbv_{n+1} - \shv_{n+1}) \Vert_{L^2(\Sigma^\star)} = O\big(\frac{1}{n^{3/2}\log^2 n}\big)$,
 \item $\Vert (\smu_n - \shmu_n) - (\smu_{n+1} - \shmu_{n+1}) \Vert_{L^2(\Sigma^\star)} = O\big(\frac{1}{n^{3/2}\log^2 n}\big)$.
\end{enumerate}
\end{Proposition}

\begin{proof}
The proof is for the most part taken from \cite[Proposition~B.2]{CD}, with the only difference being the contribution from $\Sigl =\partial U$ instead of $\Sigma \cap U$.

For (i), let us write
\begin{align}\label{IntDecomp}
 \Vert \shmu_n (\sbv_n - \shv_n) \Vert_{L^2(\Sigma^\star)}^2 = \Vert \shmu_n (\sbv_n - \shv_n) \Vert_{L^2(\Sigl)}^2 + \Vert \shmu_n (\sbv_n - \shv_n) \Vert_{L^2(\Sigr)}^2.
\end{align}
As $\hmu_n|_{s\in \Sigr} = \shmu_n|_{s\in\Sigr}$, $\bv_n|_{s\in \Sigr} = \sbv_n|_{s\in\Sigr}$ and $\hv_n|_{s\in \Sigr} = \shv_n|_{s\in\Sigr}$, one can conclude from \cite[equation~5.1]{CD} that
\begin{align} \label{Estimate05}
 \Vert \shmu_n (\sbv_n - \shv_n) \Vert_{L^2(\Sigr)} = O\left(\frac{1}{n^{1/2}\log^2 n}\right).
\end{align}
Hence it remains to consider the term $\Vert \shmu_n (\sbv_n - \shv_n) \Vert_{L^2(\Sigl)}$. Recall that for $s \in \Sigl$ we have $\sbv_n(s) = \bP^{(n)}(s)$ and $\shv_n(s) = \hP^{(n)}(s)$ and so it follows from Lemma~\ref{PropP}\,(ii) that
\begin{align}\label{jumpsComparison}
 \Vert \sbv_n - \shv_n \Vert_{L^\infty(\Sigl)} = O\big(n^{-1}\big).
\end{align}
Moreover, from Lemma~\ref{PropP}\,(iii) it also follows that $\shmu_n(s) = \hQ^{(n)}(s) \big[\hP^{(n)}(s)\big]^{-1}$ is uniformly bounded for $s \in \Sigl$. Hence we conclude
\begin{align*}
 \Vert \shmu_n (\sbv_n - \shv_n) \Vert_{L^2(\Sigl)} = O\bigl(n^{-1}\bigr)
\end{align*}
which together with \eqref{Estimate05} proves (i).

Point (ii) follows from (i) by considering
\begin{align*}
 \smu_n - \shmu_n &= \big(1-\sCauv^{(n)}\big)^{-1} \Id - \big(1-\sCauhv^{(n)}\big)^{-1}\Id
= \big(1-\sCauv^{(n)}\big)^{-1}\big(\sCauv^{(n)}-\sCauhv^{(n)}\big) \big(1-\sCauhv^{(n)}\big)^{-1}\Id
 \\
 &= \big(1-\sCauv^{(n)}\big)^{-1}\big(\sCauv^{(n)}-\sCauhv^{(n)}\big)\shmu_n= \big(1-\sCauv^{(n)}\big)^{-1}\sCau^-(\shmu_n(\sbv_n-\shv_n)).
\end{align*}
As $\big(1-\sCauv^{(n)}\big)^{-1}$, $\sCau^-$ are bounded operators (uniformly in $n$), it follows that{\samepage
\begin{align*}
 \Vert \smu_n - \shmu_n \Vert_{L^2(\Sigma^\star)} \lesssim \Vert \shmu_n(\sbv_n-\shv_n) \Vert_{L^2(\Sigma^\star)} = O\left(\frac{1}{n^{1/2}\log^2 n}\right)
\end{align*}
showing (ii).}

For (iii), we will again decompose the norm as in \eqref{IntDecomp}. As before, following the arguments found in \cite[Proposition~B.2]{CD}, we conclude that
\begin{align*}
\Vert \shmu_n (\sbv_n - \shv_n) - \shmu_{n+1} (\sbv_{n+1} - \shv_{n+1}) \Vert_{L^2(\Sigr)} = O\left(\frac{1}{n^{3/2}\log^2 n}\right).
\end{align*}
For the remaining term, we will use the fact that the local parametrices $\bP^{(n)}(s)$, $\hP^{(n)}(s)$ have an infinite series expansion in powers of $n^{-1}$ which is uniform on $\Sigl$, see \cite[equation~(8.2)]{KMVV}:
\begin{align*}
& v_n^\star(s)= P^{(n)}(s) \sim \left( \Id + \sum_{k=1}^\infty \frac{\Delta_k(s)}{n^k}\right) N(s),
 \\
& \shv_n(s)= \hP^{(n)}(s) \sim \left( \Id + \sum_{k=1}^\infty \frac{\widehat \Delta_k(s)}{n^k}\right) N(s),
\end{align*}
where $\Delta_k$ and $\widehat \Delta_k$ are $n$ independent and can be explicitly computed. In particular, \begin{align}\label{vinftyEst}
 \sbv_n(s) - \sbv_{n+1}(s) = O\big(n^{-2}\big), \qquad \shv_n(s) - \shv_{n+1}(s) = O\big(n^{-2}\big)
\end{align}
uniformly for $s \in \Sigl$. Hence, using the fact $\shmu_n(s) = \Id + O\big(n^{-1}\big)$ uniformly for $s \in \Sigl$ (see~\eqref{ParaConv}), we can write
\begin{align*}
 \Vert \shmu_n (\sbv_n - \shv_n) - \shmu_{n+1} (\sbv_{n+1} - \shv_{n+1}) \Vert_{L^2(\Sigl)}
 \lesssim \Vert (\shmu_n-\shmu_{n+1}) (\sbv_n-\shv_n) \Vert_{L^2(\Sigl)} + O\big(n^{-2}\big).
\end{align*}
Additionally, it follows that $\shmu_n(s)-\shmu_{n+1}(s) = O\big(n^{-1}\big)$, (even $O\big(n^{-2}\big)$ see \cite[equation~(8.7)]{KMVV}) uniformly for $s \in \Sigl$. Together with \eqref{jumpsComparison}, we can conclude that
\begin{align*}
 \Vert (\shmu_n-\shmu_{n+1}) (\sbv_n-\shv_n) \Vert_{L^2(\Sigl)} = O\big(n^{-2}\big),
\end{align*}
which proves (iii).

In order to prove (iv), we first write
\begin{align*}
 (\smu_n - \shmu_n)-(\smu_{n+1} - \shmu_{n+1})={}&\big(\big(1-\sCauv^{(n)}\big)^{-1} \Id - \big(1-\sCauhv^{(n)}\big)^{-1}\Id\big)
 \\
 &-\big(\big(1-\sCauv^{(n+1)}\big)^{-1} \Id - \big(1-\sCauhv^{(n+1)}\big)^{-1}\Id\big)
 \\
={}& \big(1-\sCauv^{(n)}\big)^{-1}\sCau^-(\shmu_n(\sbv_n-\shv_n))
 \\
 & -\big(1-\sCauv^{(n+1)}\big)^{-1}\sCau^-(\shmu_{n+1}(\sbv_{n+1}-\shv_{n+1}))
 \\
={}& \big(\big(1-\sCauv^{(n)}\big)^{-1}-\big(1-\sCauv^{(n+1)}\big)^{-1}\big)\sCau^-(\shmu_n(\sbv_n-\shv_n))
 \\
 & + \big(1-\sCauv^{(n+1)}\big)^{-1}\sCau^-(\shmu_n(\sbv_n-\shv_n)-\shmu_{n+1}(\sbv_{n+1}-\shv_{n+1})).
\end{align*}
From the uniform boundedness of $\big(1-\sCauv^{(n+1)}\big)^{-1}$, $\sCau^-$ and point (iii), it follows that
\begin{align*}
 \big\Vert \big(1-\sCauv^{(n+1)}\big)^{-1}\sCau^-(\shmu_n(\sbv_n-\shv_n)-\shmu_{n+1}(\sbv_{n+1}-\shv_{n+1})) \big\Vert_{L^2(\Sigma^\star)} = O\left(\frac{1}{n^{3/2}\log^2 n}\right).
\end{align*}
For the remaining term, we have
\begin{align}
 &\big(\big(1-\sCauv^{(n)}\big)^{-1}-\big(1-\sCauv^{(n+1)}\big)^{-1}\big)\sCau^-(\shmu_n(\sbv_n-\shv_n))\nonumber
 \\
 &\qquad= \big(\big(1-\sCauv^{(n)}\big)^{-1}\big(\sCauv^{(n)}-\sCauv^{(n+1)}\big)\big(1-\sCauv^{(n+1)}\big)^{-1}\big)\sCau^-(\shmu_n(\sbv_n-\shv_n)).\label{CC}
\end{align}
Now observe that
\begin{align}\label{CauMatrixIneq}
 \big\Vert \sCauv^{(n)}-\sCauv^{(n+1)} \big\Vert_{L^2(\Sigma^\star)\to L^2(\Sigma^\star)} \lesssim \Vert \sbv_n - \sbv_{n+1} \Vert_{L^\infty(\Sigma^\star)}.
\end{align}
It has been shown in \cite[p.~54]{CD} that $\Vert \sbv_n - \sbv_{n+1} \Vert_{L^\infty(\Sigr)}=O\big(n^{-1}\big)$, hence with \eqref{vinftyEst} it follows that
\begin{align*}
 \Vert \sbv_n - \sbv_{n+1} \Vert_{L^\infty(\Sigma^\star)} = O\big(n^{-1}\big),
\end{align*}
which implies using the bound \eqref{CauMatrixIneq}
\begin{align*}
 \big\Vert \sCauv^{(n)}-\sCauv^{(n+1)} \big\Vert_{L^2(\Sigma^\star)\to L^2(\Sigma^\star)} = O\big(n^{-1}\big).
\end{align*}
Plugging this estimate together with (i) into \eqref{CC}, we conclude that
\begin{align*}
 \big\Vert \big(\big(1-\sCauv^{(n)}\big)^{-1}-\big(1-\sCauv^{(n+1)}\big)^{-1}\big)\sCau^-(\shmu_n(\sbv_n-\shv_n))\big\Vert_{L^2(\Sigma^\star)} = O\left(\frac{1}{n^{3/2}\log^2 n}\right),
\end{align*}
which implies (iv) and finishes the proof.
\end{proof}

\begin{thebibliography}{99}
\bibitem{CD}
Conway T.O., Deift P., Asymptotics of polynomials orthogonal with respect to a
 logarithmic weight, \href{https://doi.org/10.3842/SIGMA.2018.056}{\textit{SIGMA}} \textbf{14} (2018), 056, 66~pages,
 \href{https://arxiv.org/abs/1711.01590}{arXiv:1711.01590}.


\bibitem{KMVV}
Kuijlaars A.B.J., McLaughlin K.T.-R., Van~Assche W., Vanlessen M., The
 {R}iemann--{H}ilbert approach to strong asymptotics for orthogonal
 polynomials on {$[-1,1]$}, \href{https://doi.org/10.1016/j.aim.2003.08.015}{\textit{Adv. Math.}} \textbf{188} (2004), 337--398,
 \href{https://arxiv.org/abs/math.CA/0111252}{arXiv:math.CA/0111252}.
\end{thebibliography}
\end{document}
